\documentclass[12pt,a4paper]{ctexart}
\usepackage{geometry,graphicx,amsmath,amssymb,booktabs,xcolor,hyperref,bookmark}
\usepackage{tikz}
\usetikzlibrary{shapes.geometric, arrows.meta, positioning, calc}

\geometry{left=2.5cm,right=2.5cm,top=2.5cm,bottom=2.5cm}

\title{强化学习在飞行控制中的应用：综述}
\author{}
\date{}

\begin{document}
\maketitle
\tableofcontents
\newpage

% 第1章 引言
\section{引言}
\subsection{研究背景与意义}
\subsection{问题定义与挑战}
\subsection{研究现状概述}
\subsection{本文组织结构}

% 第2章 基础理论
\section{强化学习与飞行控制基础}
\subsection{强化学习基础}
\subsection{飞行器动力学建模}
\subsection{仿真环境与基准测试}

% 第3章 多旋翼RL控制
\section{多旋翼无人机强化学习控制}
\subsection{姿态控制}
\subsection{位置与轨迹跟踪}
\subsection{避障与路径规划}
\subsection{故障容错控制}

% 第4章 固定翼RL控制
\section{固定翼飞行器强化学习控制}
\subsection{巡航与机动控制}
\subsection{失速恢复与边界保护}
\subsection{路径跟踪与制导}
\subsection{对抗不确定性的鲁棒控制}

% 第5章 VTOL RL控制
\section{垂直起降飞行器强化学习控制}
\subsection{过渡飞行控制}
\subsection{倾转旋翼控制}
\subsection{尾座式无人机控制}
\subsection{多模态切换策略}

% 第6章 实验与应用
\section{实验验证与应用}
\subsection{仿真平台与工具}
\subsection{Sim-to-Real迁移}
\subsection{实飞验证案例}
\subsection{性能对比分析}

% 第7章 未来方向
\section{未来研究方向}
\subsection{技术挑战}
\subsection{研究趋势}
\subsection{展望}

\begin{thebibliography}{99}
% 文献将在此添加
\end{thebibliography}

\end{document}
